\section{Related Work}
\label{sec:related_work}

\subsection{Reinforcement Learning for Knowledge Graphs}
MINERVA by Das et al.~\cite{ref_das2018} learns policies for multi-hop query answering. DeepPath by Xiong et al.~\cite{ref_xiong2017} uses reinforcement learning to discover reasoning paths. Lin et al.~\cite{ref_lin2018} study reward shaping for multi-hop reasoning. These methods select graph traversal actions to answer queries. Our action space instead changes graph records and acquires validators that enable later repairs.

The scheduling question is whether a learned value function can allocate a limited budget better than a fixed schedule or an observation-based heuristic. We compare these alternatives on the same executable transitions and use retrained ablations to test the roles of observed features and action masking. The controlled registry lets us reproduce each transition, while the text-derived rule study examines the candidate source separately.

\subsection{Language Models for Knowledge Engineering}
KG-BERT~\cite{ref_yao2019} uses a language model for triple classification, and KEPLER~\cite{ref_wang2021kepler} combines language modeling with knowledge embeddings. Pan et al.~\cite{ref_pan2024} review ways to connect language models and knowledge graphs. These studies motivate using text as a source of information for graph construction and completion.

Chain-of-thought prompting~\cite{ref_wei2023} and ReAct~\cite{ref_yao2023} explore structured reasoning and tool use. Our generation module uses two single-call elicitation tasks with a common output schema: deletion asks which constraints are exposed by missing spans, while augmentation asks which constraints are suggested by supplementary clauses. We evaluate the resulting candidates through archived provenance and executable patterns, rather than treating a reasoning response as a validated rule.

\subsection{Constraint Languages, Rule Mining, and Repair}
SHACL~\cite{ref_knublauch2017} and ShEx~\cite{ref_prud2014} express graph constraints. Their detection coverage depends on the shapes provided. A limited set of schema checks should therefore be distinguished from the expressiveness of the constraint language itself.

AMIE~\cite{ref_galarraga2013} mines Horn rules from graph facts, RuDiK~\cite{ref_ortona2018} studies rule discovery in knowledge bases, and Neural LP~\cite{ref_yang2017} learns logical rules through differentiable operations. Our candidates originate from source documents and include type patterns and procedural statements. Their different representation requires a validation and compilation step before they can act as graph validators.

Lin et al.~\cite{ref_lin2025} systematically evaluate LLM-based graph repair. This makes controlled constraint and repair comparisons directly relevant to our study. We focus on the order of validator acquisition and graph edits, with an additional provenance-based audit of generated candidates. An end-to-end comparison with external repair and rule-mining methods requires aligning executable constraints, input evidence, and evaluation budgets.
